% ============================================================
% Section 155: 周期性 $\delta$ 势垒的完全反射条件
% ============================================================
\section{量子力学：周期性 $\delta$ 势垒的完全反射条件}
\label{sec:periodic-delta-reflection}

\begin{modelbox}[题目：周期性 $\delta$ 势垒的完全反射]
质量为 $m$ 的粒子束以动量 $p=\hbar k$ 从 $x=-\infty$ 入射，碰到周期性 $\delta$ 势垒
\[
V(x)=\gamma\sum_{n=0}^{\infty}\delta(x-na),\quad a>0,
\]
求发生完全反射的动量值。
\end{modelbox}

\begin{tipbox}[考点快速分析]
\begin{itemize}
    \item \textbf{分区波函数}：各区为左右行波叠加，$N$ 个势垒分 $N+1$ 区
    \item \textbf{边界条件}：$\delta$ 势处波函数连续、导数跃变 $2m\gamma/\hbar^2$
    \item \textbf{完全反射条件}：$e^{2ika}=1$，相邻势垒反射波同相叠加
    \item \textbf{物理对应}：Kronig-Penney 模型禁带中心、Bragg 反射条件
\end{itemize}
\end{tipbox}

\begin{derivationbox}[标准解题过程]
\textbf{1. 有限 $N$ 个势垒的模型}

设 $N$ 个 $\delta$ 势垒位于 $x=0,a,\ldots,(N-1)a$。第 $n$ 区波函数：
\[
\psi_n(x)=D_ne^{ikx}+R_ne^{-ikx}.
\]
边界条件：$D_0=1$（入射波），$R_N=0$（最右无反射波）。

\textbf{2. 递推关系}

在 $x=na$ 处：
\begin{align*}
\psi_n(na)&=\psi_{n+1}(na),\\
\psi_{n+1}'(na)-\psi_n'(na)&=\frac{2m\gamma}{\hbar^2}\psi_n(na).
\end{align*}

定义无量纲参数 $\alpha=\dfrac{m\gamma}{\hbar^2k}$，得递推关系：
\[
R_{n+1}-R_n=(D_n-D_{n+1})e^{2inka},\quad R_{n+1}-R_n=i\alpha(R_n+D_ne^{2inka}).
\]

\textbf{3. 完全反射条件分析}

当 $e^{2ika}=1$ 时，递推关系简化。对 $n=0$ 到 $N-1$ 求和得 $D_0+R_0=D_N$。进一步递推得 $-R_0=iN\alpha D_N$。

概率流守恒：$|D_0|^2=|R_0|^2+|D_N|^2$。代入 $R_0=-iN\alpha D_N$ 得
\[
|D_N|^2=\frac{1}{1+N^2\alpha^2}.
\]
当 $N\to\infty$ 时 $|D_N|^2\to0$，透射趋于零，反射趋于 1。同时 $D_0+R_0=D_N\to0$ 得 $R_0\to-D_0$，即全反射且反射波与入射波相位差 $\pi$。

\textbf{4. 全反射的动量条件}

由 $e^{2ika}=1$ 得 $2ka=2m\pi$（$m\in\mathbb{Z}$），即
\begin{equation}
\boxed{k=\frac{m\pi}{a},\quad p=\frac{m\pi\hbar}{a}\quad(m=1,2,3,\ldots).}
\end{equation}

对应的波长 $\lambda=2a/m$，相邻势垒间距恰好是半波长的整数倍，各势垒反射波同相叠加，发生相长干涉，形成完全反射。

\textbf{物理图像}：这实际上就是 Kronig-Penney 模型中能带间隙（禁带）中心的动量值，对应 Bloch 波在布里渊区边界发生 Bragg 反射的条件。
\end{derivationbox}

\begin{formulabox}[答案汇总]
\begin{center}
\begin{tabular}{ll}
\toprule
\textbf{结论} & \textbf{结果} \\
\midrule
完全反射条件 & $e^{2ika}=1$ \\
动量取值 & $p=m\pi\hbar/a$（$m=1,2,3,\ldots$） \\
波长条件 & $\lambda=2a/m$（Bragg 条件） \\
\bottomrule
\end{tabular}
\end{center}
\end{formulabox}

\begin{insightbox}[物理图像]
\begin{itemize}
    \item 完全反射条件 $e^{2ika}=1$ 是使所有 $\delta$ 势垒的反射波同相叠加的充分条件。在实际晶体中，这对应 X 射线 Bragg 衍射的量子力学类比。
    \item Kronig-Penney 模型中的能带间隙正是源于这种周期性势场的 Bragg 反射。在禁带中心，电子波无法传播，被完全反射。
    \item 该结果也是理解光子晶体、声子晶体等人工周期结构中``带隙''现象的量子力学基础。
\end{itemize}
\end{insightbox}
